\documentclass[11pt, a4paper]{article}

% ─── Encodage et langue ───
\usepackage[utf8]{inputenc}
\usepackage[T1]{fontenc}
\usepackage[french]{babel}

% ─── Maths ───
\usepackage{amsmath, amssymb, amsthm}

% ─── Mise en page ───
\usepackage[margin=2.5cm]{geometry}
\usepackage{enumitem}
\usepackage{booktabs}
\usepackage{multirow}
\usepackage{graphicx}
\usepackage{float}
\usepackage{caption}
\usepackage{subcaption}
\usepackage{xcolor}
\usepackage{hyperref}
\hypersetup{colorlinks=true, linkcolor=blue!60!black, citecolor=green!50!black, urlcolor=blue!70!black}
% ─── Code R ───
\usepackage{listings}
\lstset{
  language=R,
  basicstyle=\small\ttfamily,
  keywordstyle=\color{blue!70!black}\bfseries,
  commentstyle=\color{green!50!black}\itshape,
  stringstyle=\color{red!60!black},
  numbers=left, numberstyle=\tiny\color{gray},
  frame=single, framerule=0.4pt,
  breaklines=true, showstringspaces=false,
  tabsize=2
}

% ─── Commandes personnalisées ───
\newcommand{\R}{\mathbb{R}}
\newcommand{\E}{\mathbb{E}}
\newcommand{\Var}{\mathrm{Var}}
\newcommand{\Ltwo}{L^2(\mathcal{T})}
\newcommand{\Dp}{D_p}
\newcommand{\Dw}{D_w}
\newcommand{\DK}{D_K}

% ─── Théorèmes ───
\theoremstyle{definition}
\newtheorem{definition}{Définition}[section]
\theoremstyle{remark}
\newtheorem{remarque}{Remarque}[section]

% ============================================================================
\title{%
  \textbf{Rapport complet} \\[0.3em]
  \Large Classification non supervisée de données fonctionnelles mixtes \\[0.2em]
  \large Application au jeu de données Canadian Weather
}
\author{Pierre}
\date{Mars 2026}

\begin{document}
\maketitle

\begin{abstract}
Ce rapport documente l'intégralité du travail réalisé sur le jeu de données
\texttt{CanadianWeather} du package R \texttt{fda}. L'objectif est de tester et
comparer trois stratégies de classification non supervisée pour des données
\textbf{mixtes} : une composante \textbf{fonctionnelle} (courbes de température
journalière) et une composante \textbf{vectorielle} (latitude, longitude,
précipitations). Le pipeline complet va du lissage par B-splines pénalisées
jusqu'au clustering, en passant par l'ACP fonctionnelle et le calcul de
plusieurs distances. Les résultats sont évalués par silhouette et ARI (Adjusted
Rand Index) en référence aux 4 régions climatiques connues. Ce travail sert
de \textbf{prototype} pour l'application future aux données océanographiques
du SHOM.
\end{abstract}

\tableofcontents
\newpage

% ============================================================================
% ============================================================================
\part{Le jeu de données}
% ============================================================================
% ============================================================================

\section{Présentation de Canadian Weather}

Le jeu de données \texttt{CanadianWeather}, inclus dans le package R
\texttt{fda} \cite{ramsay2005}, contient les moyennes climatiques journalières
pour $n = 35$ stations météorologiques réparties sur l'ensemble du territoire
canadien.

\subsection{Composante fonctionnelle : températures}

Pour chaque station $i \in \{1, \ldots, 35\}$, on dispose de la température
moyenne quotidienne sur un an :
\[
  Y_{ij} = X_i(t_j) + \varepsilon_{ij}, \quad j = 1, \ldots, 365,
\]
où $t_j$ est le $j$-ème jour de l'année. Cela produit une matrice
$\mathbf{Y} \in \R^{365 \times 35}$. Chaque colonne est une
\textbf{courbe discrète} de température.

\subsection{Composante vectorielle : variables géographiques}

Pour chaque station, on dispose de trois variables scalaires :
\[
  Z_i = \begin{pmatrix} \text{latitude}_i \\ \text{longitude}_i \\ \text{précipitations}_i \end{pmatrix} \in \R^3.
\]
Les précipitations sont calculées comme la moyenne journalière
($\bar{\text{Precip}}_i = \frac{1}{365} \sum_j \text{Precip}_{ij}$).
Ces variables forment la matrice $\mathbf{Z} \in \R^{35 \times 3}$.

\subsection{Vérité terrain : les 4 régions}

Chaque station est classée dans l'une des 4 régions climatiques :
\begin{center}
\begin{tabular}{lrl}
  \toprule
  \textbf{Région} & \textbf{Nb stations} & \textbf{Caractéristiques} \\
  \midrule
  Arctic       & 3  & Nord, froid extrême, faible amplitude \\
  Atlantic     & 15 & Côte est, hivers modérés, précipitations élevées \\
  Continental  & 12 & Centre, hivers très froids, étés chauds \\
  Pacific      & 5  & Côte ouest, climat tempéré, faible amplitude \\
  \bottomrule
\end{tabular}
\end{center}
Ces étiquettes ne sont \textbf{jamais utilisées} par les algorithmes de
clustering (classification non supervisée). Elles servent uniquement
\textbf{a posteriori} pour évaluer la qualité des regroupements trouvés,
via l'ARI (section~\ref{sec:ari}).

\subsection{Pourquoi ce jeu de données ?}

Canadian Weather est le jeu de données canonique en analyse de données
fonctionnelles (FDA). Il présente plusieurs avantages pour notre travail :
\begin{itemize}[nosep]
  \item \textbf{Petite taille} ($n = 35$) : permet des itérations rapides.
  \item \textbf{Données mixtes} : courbes + variables scalaires.
  \item \textbf{Vérité terrain connue} : les 4 régions climatiques.
  \item \textbf{Structure claire} : les courbes arctiques, continentales,
    atlantiques et pacifiques sont visuellement distinctes.
\end{itemize}

\newpage
% ============================================================================
% ============================================================================
\part{Pipeline complet}
% ============================================================================
% ============================================================================

Le pipeline est organisé en 6 scripts exécutés séquentiellement :
\begin{center}
\begin{tabular}{rll}
  \toprule
  \textbf{Script} & \textbf{Étape} & \textbf{Sortie principale} \\
  \midrule
  \texttt{00\_preprocess.R} & Chargement des données & $\mathbf{Y}$, $\mathbf{Z}$, régions \\
  \texttt{01\_lissage.R}    & Lissage B-splines      & $\hat{X}_i(t) \in \Ltwo$ \\
  \texttt{02\_fpca.R}       & FPCA                   & Scores $\xi_{ik}$ \\
  \texttt{03\_distances.R}  & Calcul des distances   & $D_0, D_1, \Dp, D_s, \Dw, \DK$ \\
  \texttt{04\_clustering.R} & Clustering (3 strat.)  & Labels, ARI, silhouette \\
  \texttt{05\_visualisation.R} & Figures              & 8 figures PNG \\
  \bottomrule
\end{tabular}
\end{center}

% ============================================================================
\section{Étape 0 : Prétraitement}
\label{sec:preprocess}

Le script \texttt{00\_preprocess.R} charge le dataset et prépare les objets :

\begin{itemize}[nosep]
  \item \texttt{Y\_brut} $\in \R^{365 \times 35}$ : températures brutes.
  \item \texttt{Z} $\in \R^{35 \times 3}$ : latitude, longitude, précipitations.
  \item \texttt{regions} : facteur à 4 niveaux (Arctic, Atlantic, Continental, Pacific).
  \item \texttt{noms\_stations} : noms des 35 stations.
\end{itemize}

Les variables $Z$ sont aussi \textbf{standardisées} (\texttt{Z\_scaled = scale(Z)})
pour que latitude, longitude et précipitations soient sur une même échelle
avant le calcul des distances vectorielles.

% ============================================================================
\section{Étape 1 : Lissage par B-splines pénalisées}
\label{sec:lissage}

\subsection{Objectif}

Passer des observations discrètes $Y_{ij}$ à des courbes continues
$\hat{X}_i(t)$ évaluables en tout point $t \in [1, 365]$.

\subsection{Modèle}

On décompose chaque courbe sur une base de B-splines :
\[
  \hat{X}_i(t) = \sum_{k=1}^{d} \hat{\theta}_{ik} \, \psi_k(t),
\]
où les $\psi_k$ sont des B-splines d'ordre 4 (cubiques). Les coefficients
$\hat{\theta}_{ik}$ sont estimés par moindres carrés pénalisés :
\begin{equation}
  \hat{\boldsymbol{\theta}}_i = \arg\min_{\boldsymbol{\theta}}
  \left\{
    \sum_{j=1}^{365} \left( Y_{ij} - \sum_{k=1}^{d} \theta_k \psi_k(t_j) \right)^2
    + \lambda \int_1^{365} \left[ \hat{X}_i''(t) \right]^2 dt
  \right\}.
  \label{eq:lissage}
\end{equation}

\subsection{Choix des hyperparamètres}

\paragraph{Nombre de bases : $d = 65$.}
On place environ une base par semaine (365/65 $\approx$ 5.6 jours). Ce nombre
est volontairement \textbf{sur-dimensionné} : la base a la capacité de produire
des oscillations rapides, mais le paramètre $\lambda$ régule la régularité.
C'est la philosophie \emph{<<~beaucoup de bases + pénalité~>>} recommandée par
Ramsay \& Silverman \cite{ramsay2005}.

\paragraph{Paramètre de pénalité : $\lambda$ par GCV.}
Le paramètre $\lambda$ est choisi par \textbf{validation croisée généralisée}
(GCV). On teste une grille logarithmique de 80 valeurs entre $10^{-2}$ et
$10^{8}$, et on minimise :
\[
  \text{GCV}(\lambda) = \frac{1}{n} \sum_{i=1}^{n}
  \frac{\| Y_i - \hat{X}_i^\lambda \|^2}{(1 - \text{df}(\lambda)/N)^2},
\]
où $\text{df}(\lambda) = \text{tr}(S_\lambda)$ est le nombre de degrés de
liberté effectifs. Le minimum est atteint en :

\begin{center}
\fbox{$\lambda_{\text{GCV}} = 8.15 \quad (\log_{10} \lambda = 0.91)$}
\end{center}

\begin{remarque}
La valeur $\lambda = 8.15$ est \textbf{faible}, ce qui signifie que les courbes
de température canadiennes sont relativement lisses par nature et nécessitent
peu de régularisation. Le minimum GCV n'est pas au bord de la grille, ce qui
confirme la fiabilité du choix.
\end{remarque}

\subsection{Résultat}

Le lissage produit un objet \texttt{fd} (functional data) contenant 35 courbes
$\hat{X}_i(t)$ dans $\Ltwo$, dérivables à tout ordre.

% ============================================================================
\section{Étape 2 : ACP fonctionnelle (FPCA)}
\label{sec:fpca}

\subsection{Principe}

L'ACP fonctionnelle diagonalise l'opérateur de covariance :
\[
  \int_{\mathcal{T}} C(s,t) \, \varphi_k(s) \, ds = \lambda_k \, \varphi_k(t),
\]
où $C(s,t) = \frac{1}{n} \sum_i [X_i(s) - \mu(s)][X_i(t) - \mu(t)]$ est la
surface de covariance.

\subsection{Choix de $K$}

On retient les $K$ premières composantes telles que la variance cumulée atteigne
$\geq 95\%$ :

\begin{center}
\begin{tabular}{crr}
  \toprule
  \textbf{PC} & \textbf{Var. expliquée} & \textbf{Var. cumulée} \\
  \midrule
  $\varphi_1$ & 88.5\% & 88.5\% \\
  $\varphi_2$ &  8.5\% & 97.0\% \\
  $\varphi_3$ &  2.0\% & 99.0\% \\
  $\varphi_4$ &  0.5\% & 99.5\% \\
  \bottomrule
\end{tabular}
\end{center}

\begin{center}
\fbox{$K = 2$ composantes retenues (variance cumulée $= 97.0\%$)}
\end{center}

\subsection{Interprétation des fonctions propres}

\paragraph{$\varphi_1(t)$ — Mode de niveau (88.5\%).}
Cette fonction propre est approximativement constante et positive sur tout
l'intervalle. Elle capture le \textbf{niveau moyen de température} : une station
avec un score $\xi_{i1}$ élevé est globalement plus chaude.

\paragraph{$\varphi_2(t)$ — Mode de saisonnalité (8.5\%).}
Cette fonction est positive en hiver et négative en été. Elle capture le
\textbf{contraste saisonnier} : un score $\xi_{i2}$ élevé correspond à un
écart hivernal plus marqué (climat continental), tandis qu'un score faible
indique un climat plus tempéré (maritime/pacifique).

\subsection{Scores}

Chaque station est résumée par un vecteur de 2 scores :
\[
  \xi_i = (\xi_{i1}, \xi_{i2}) \in \R^2, \quad
  \xi_{ik} = \int_1^{365} [\hat{X}_i(t) - \mu(t)] \, \varphi_k(t) \, dt.
\]
On passe de 35 courbes dans $\Ltwo$ (dimension infinie) à 35 points dans
$\R^2$ (dimension finie). Cette réduction est \textbf{optimale} au sens de la
variance capturée (théorème de Karhunen-Loève).

% ============================================================================
\section{Étape 3 : Calcul des distances}
\label{sec:distances}

L'architecture des distances est organisée en \textbf{deux niveaux} :

\begin{figure}[H]
\centering
\begin{tabular}{c}
\fbox{\parbox{0.85\textwidth}{
\textbf{Niveau 1 — Distance fonctionnelle $\Dp(\alpha)$} \\[0.3em]
$D_0$ (L² entre courbes, niveau) \\
$D_1$ (L² entre dérivées, forme) \\
$\Dp(\alpha) = \sqrt{(1-\alpha) \cdot \tilde{D}_0^2 + \alpha \cdot \tilde{D}_1^2}$ \\[0.5em]
\textbf{Niveau 2 — Combinaison fonctionnel + vectoriel} \\[0.3em]
$D_s$ (euclidienne sur $Z$ standardisé) \\
$\Dw(\alpha, \omega) = \sqrt{\omega \cdot \Dp(\alpha)^2 + (1-\omega) \cdot \tilde{D}_s^2}$ \\
$\DK(\alpha) = \sqrt{K(i,i) + K(j,j) - 2K(i,j)}$ avec $K_f = \exp(-\Dp(\alpha)^2 / 2\sigma_f^2)$
}}
\end{tabular}
\caption{Architecture à deux niveaux des distances.}
\label{fig:archi_distances}
\end{figure}

\subsection{Distance $D_0$ : L² entre courbes}

\begin{definition}[Distance $D_0$]
\[
  D_0(i,j) = \sqrt{\int_1^{365} [\hat{X}_i(t) - \hat{X}_j(t)]^2 \, dt}
\]
\end{definition}

C'est la norme $L^2$ de la différence. Elle mesure la proximité
\textbf{globale} des profils. En pratique, l'intégrale est approchée par une
somme de Riemann sur 1000 points équidistants.

\textbf{Résultat :} $D_0 \in [3.9, \, 532.1]$.

\subsection{Distance $D_1$ : L² entre dérivées}

\begin{definition}[Distance $D_1$]
\[
  D_1(i,j) = \sqrt{\int_1^{365} [\hat{X}_i'(t) - \hat{X}_j'(t)]^2 \, dt}
\]
\end{definition}

Capture les différences de \textbf{dynamique} (vitesse de changement).
Deux courbes au même niveau mais dont l'une chute plus vite en automne
seront proches en $D_0$ mais éloignées en $D_1$.

\textbf{Résultat :} $D_1 \in [0.51, \, 6.61]$.

\begin{remarque}
Les échelles de $D_0$ et $D_1$ sont très différentes (facteur $\sim 80$).
C'est pourquoi une \textbf{normalisation} est indispensable avant de les
combiner.
\end{remarque}

\subsection{Distance $\Dp(\alpha)$ : combinaison fonctionnelle}

\begin{definition}[Distance $\Dp$]
Soient $\tilde{D}_0 = D_0 / \max(D_0)$ et $\tilde{D}_1 = D_1 / \max(D_1)$ les
distances normalisées par leur maximum. Alors :
\[
  \Dp(\alpha) = \sqrt{(1-\alpha) \cdot \tilde{D}_0^2 + \alpha \cdot \tilde{D}_1^2},
  \qquad \alpha \in [0, 1].
\]
\end{definition}

Le paramètre $\alpha$ contrôle le \textbf{compromis niveau/forme} :
\begin{itemize}[nosep]
  \item $\alpha = 0$ : $\Dp = \tilde{D}_0$ (seul le niveau compte),
  \item $\alpha = 1$ : $\Dp = \tilde{D}_1$ (seule la forme compte),
  \item $\alpha \in (0,1)$ : compromis.
\end{itemize}

\subsection{Distance $D_s$ : euclidienne sur $Z$}

\begin{definition}[Distance $D_s$]
\[
  D_s(i,j) = \| Z_i^{\text{std}} - Z_j^{\text{std}} \|_2,
\]
où $Z_i^{\text{std}} = (Z_i - \bar{Z}) / \sigma_Z$ est le vecteur standardisé.
\end{definition}

\textbf{Résultat :} $D_s \in [0.07, \, 5.68]$.

\subsection{Distance $\Dw(\alpha, \omega)$ : mixte pondérée (stratégie B)}

\begin{definition}[Distance $\Dw$]
Soit $\tilde{D}_s = D_s / \max(D_s)$. Alors :
\[
  \Dw(\alpha, \omega) = \sqrt{\omega \cdot \Dp(\alpha)^2 + (1-\omega) \cdot \tilde{D}_s^2},
  \qquad \alpha, \omega \in [0, 1].
\]
\end{definition}

C'est un système à \textbf{deux paramètres} :
\begin{itemize}[nosep]
  \item $\alpha$ : mix D0/D1 \textbf{dans} le fonctionnel (niveau vs forme),
  \item $\omega$ : mix fonctionnel/vectoriel (courbes vs variables Z).
\end{itemize}

Les paramètres sont choisis par \textbf{grid search} sur
$\alpha \times \omega \in \{0, 0.1, \ldots, 1\}^2$ (121 combinaisons),
en maximisant la silhouette.

\subsection{Distance $\DK(\alpha)$ : à noyaux (stratégie C)}

\begin{definition}[Distance $\DK$]
On définit le noyau produit :
\begin{align*}
  K_f(i,j) &= \exp\!\left(-\frac{\Dp(\alpha)_{ij}^2}{2\sigma_f^2}\right), \qquad
  \sigma_f = \text{médiane}\{\Dp(\alpha)_{ij} : i < j\}, \\
  K_s(i,j) &= \exp\!\left(-\frac{D_s(i,j)^2}{2\sigma_s^2}\right), \qquad
  \sigma_s = \text{médiane}\{D_s(i,j) : i < j\}, \\
  K(i,j) &= K_f(i,j) \cdot K_s(i,j).
\end{align*}
La distance associée est :
\[
  \DK(i,j) = \sqrt{K(i,i) + K(j,j) - 2K(i,j)}.
\]
\end{definition}

Les largeurs de bande $\sigma_f$ et $\sigma_s$ sont fixées par l'heuristique
de la médiane \cite{jaakkola1999}. Le seul paramètre libre est $\alpha$,
optimisé par grid search sur $\{0, 0.1, \ldots, 1\}$.

% ============================================================================
\section{Étape 4 : Clustering}
\label{sec:clustering}

Trois stratégies sont comparées, plus deux baselines.

\subsection{Stratégie A : FPCA + concaténation + k-means}

On concatène les scores FPCA standardisés et les variables $Z$ standardisées :
\[
  W_i = \left(\frac{\xi_{i1}}{\sigma_{\xi_1}}, \frac{\xi_{i2}}{\sigma_{\xi_2}},
  Z_i^{\text{std}}\right) \in \R^{K+3} = \R^5.
\]
On applique \textbf{k-means} avec $k = 4$ clusters et 25 initialisations
aléatoires (\texttt{nstart = 25}).

\paragraph{Pourquoi k-means ?} k-means minimise la variance intra-cluster et
opère sur des vecteurs. Il n'utilise pas de matrice de distances : il travaille
directement dans $\R^5$.

\subsection{Stratégie B : $\Dw(\alpha, \omega)$ + PAM}

On construit $\Dw(\alpha, \omega)$ pour chaque couple $(\alpha, \omega)$ dans
la grille $\{0, 0.1, \ldots, 1\}^2$, puis on applique \textbf{PAM}
(Partitioning Around Medoids) avec $k = 4$.

\paragraph{Pourquoi PAM et pas k-means ?}
k-means n'accepte que des vecteurs de $\R^p$ et utilise la distance euclidienne.
PAM est l'extension de k-means à des \textbf{matrices de distances arbitraires}.
Au lieu de calculer des centroïdes (moyennes), PAM sélectionne des
\textbf{médoïdes} (individus réels du jeu de données) qui minimisent la somme
des distances aux points de leur cluster.

\paragraph{Grid search 2D.}
Pour chaque couple $(\alpha, \omega)$, on calcule la silhouette moyenne. Le
couple optimal est celui qui maximise ce critère. Cela représente
$11 \times 11 = 121$ exécutions de PAM.

\subsection{Stratégie C : $\DK(\alpha)$ + PAM}

On construit $\DK(\alpha)$ pour chaque $\alpha \in \{0, 0.1, \ldots, 1\}$, puis
PAM avec $k = 4$. Le $\alpha$ optimal maximise la silhouette.

\subsection{Baselines}

Deux baselines isolent chaque source d'information :
\begin{itemize}[nosep]
  \item \textbf{$D_0$ seul} : PAM sur la distance fonctionnelle brute (pas de Z).
  \item \textbf{$D_s$ seul} : PAM sur la distance vectorielle (pas de courbes).
\end{itemize}

% ============================================================================
\section{Métriques d'évaluation}
\label{sec:eval}

\subsection{Silhouette moyenne}

Pour un individu $i$ assigné au cluster $C_k$ :
\[
  s(i) = \frac{b(i) - a(i)}{\max(a(i), b(i))},
\]
où $a(i) = $ distance moyenne aux individus du même cluster, et $b(i) = $
distance moyenne au cluster voisin le plus proche. La silhouette moyenne
$\bar{s} = \frac{1}{n}\sum_i s(i)$ mesure la qualité \textbf{intrinsèque}
du clustering :
\begin{itemize}[nosep]
  \item $\bar{s} \approx 1$ : clusters très bien séparés,
  \item $\bar{s} \approx 0$ : frontières floues,
  \item $\bar{s} < 0$ : individus mal classés.
\end{itemize}

\subsection{ARI (Adjusted Rand Index)}
\label{sec:ari}

L'ARI compare la partition trouvée $\hat{C}$ à la vérité terrain $C^*$
(les 4 régions). C'est une version corrigée pour le hasard de l'indice de Rand :
\[
  \text{ARI} = \frac{\text{Rand} - \E[\text{Rand}]}{\max(\text{Rand}) - \E[\text{Rand}]}.
\]
\begin{itemize}[nosep]
  \item $\text{ARI} = 1$ : accord parfait avec la vérité terrain,
  \item $\text{ARI} = 0$ : accord équivalent au hasard,
  \item $\text{ARI} < 0$ : pire que le hasard.
\end{itemize}

\begin{remarque}
L'ARI est utilisé \textbf{a posteriori} pour mesurer à quel point le clustering
non supervisé retrouve les régions connues. Il ne guide \textbf{pas} le
clustering : c'est la silhouette qui sert de critère de sélection des
paramètres $(\alpha, \omega)$. Cette distinction est fondamentale : le
clustering reste non supervisé.
\end{remarque}

\newpage
% ============================================================================
% ============================================================================
\part{Résultats}
% ============================================================================
% ============================================================================

\section{Tableau comparatif}

\begin{table}[H]
\centering
\begin{tabular}{llcc}
  \toprule
  & \textbf{Stratégie} & \textbf{Silhouette} & \textbf{ARI} \\
  \midrule
  \multirow{2}{*}{\textit{Baselines}}
  & $D_0$ seul (fonctionnel pur)  & 0.285 & 0.229 \\
  & $D_s$ seul (vectoriel pur)    & 0.448 & 0.616 \\
  \midrule
  & A : FPCA + Z $\to$ k-means   & 0.395 & \textbf{0.748} \\
  & B : $\Dw(\alpha\!=\!0, \omega\!=\!0)$ $\to$ PAM & \textbf{0.448} & 0.616 \\
  & C : $\DK(\alpha\!=\!0)$ $\to$ PAM & 0.329 & 0.686 \\
  \bottomrule
\end{tabular}
\caption{Résultats de clustering sur Canadian Weather ($k = 4$, seed = 42).}
\label{tab:resultats}
\end{table}

\section{Analyse de chaque stratégie}

\subsection{Stratégie A : meilleur ARI (0.748)}

La stratégie A obtient le meilleur ARI. Le k-means dans l'espace
$\R^5 = \R^{K+3}$ parvient à exploiter conjointement les scores FPCA
(résumé des courbes) et les variables géographiques. La silhouette est
correcte (0.395) mais inférieure à celle de $D_s$ seul.

\paragraph{Pourquoi ?} K-means minimise la variance intra-cluster en utilisant
la distance euclidienne dans l'espace concaténé. Les coordonnées
géographiques (surtout la latitude) et les scores FPCA (surtout $\xi_1$,
le niveau moyen) sont des proxies complémentaires de la région climatique.
Le k-means parvient à trouver un compromis géométrique efficace.

\subsection{Stratégie B : retombe sur $D_s$ seul ($\omega = 0$)}

Le grid search 2D sur $(\alpha, \omega)$ identifie l'optimum en
$\alpha = 0, \omega = 0$, ce qui signifie :
\begin{itemize}[nosep]
  \item $\omega = 0$ : le fonctionnel est ignoré $\Rightarrow$ $\Dw = \tilde{D}_s$,
  \item $\alpha = 0$ : $D_1$ est ignoré $\Rightarrow$ $\Dp = \tilde{D}_0$ (sans importance puisque $\omega = 0$).
\end{itemize}

\paragraph{Conclusion.} La stratégie B est identique au baseline $D_s$ seul.
La composante fonctionnelle n'apporte rien \textbf{dans ce cadre spécifique}.

\subsection{Stratégie C : ARI intermédiaire (0.686)}

Le noyau produit $K = K_f \cdot K_s$ combine fonctionnel et vectoriel sans
paramètre $\omega$. Le choix de $\alpha = 0$ signifie que le noyau fonctionnel
utilise $D_0$ (courbes brutes).

L'ARI (0.686) est supérieur à $D_s$ seul (0.616), ce qui montre que le noyau
parvient à \textbf{intégrer} les deux sources d'information de manière
non triviale, même quand $D_s$ domine.

% ============================================================================
\section{Diagnostic : pourquoi les variables vectorielles écrasent le fonctionnel}
\label{sec:critique}

Le résultat central de cette étude est que \textbf{les variables vectorielles
$Z$ écrasent la composante fonctionnelle}. Ce n'est pas un défaut de la
méthode — c'est une propriété du jeu de données Canadian Weather. Nous le
prouvons ici par cinq arguments quantitatifs indépendants.

% ─────────────────────────────────────────────────────────────────────────────
\subsection{Argument 1 : corrélation écrasante latitude–température}

La température moyenne annuelle de chaque station, $\bar{T}_i = \frac{1}{365}
\sum_j Y_{ij}$, est un résumé scalaire de la courbe~$X_i(t)$. Sa corrélation
avec les variables vectorielles est :

\begin{center}
\begin{tabular}{lc}
  \toprule
  \textbf{Variable $Z$} & $\boldsymbol{r}$ \textbf{(Pearson)} \\
  \midrule
  Latitude      & $-0.884$ \\
  Longitude     & $-0.145$ \\
  Précipitations & $+0.577$ \\
  \bottomrule
\end{tabular}
\end{center}

La corrélation latitude–température de $r = -0.884$ signifie que \textbf{78\%
de la variance} ($r^2 = 0.78$) de la température moyenne est expliquée par la
latitude seule. Autrement dit, le \textbf{contenu informatif principal} des
courbes (le niveau moyen) est déjà capturé par une seule variable scalaire.
Les courbes sont donc largement \textbf{redondantes} avec $Z$.

% ─────────────────────────────────────────────────────────────────────────────
\subsection{Argument 2 : ANOVA — les variables $Z$ séparent les régions}

On mesure, pour chaque variable vectorielle, la fraction de variance expliquée
par les 4 régions ($\eta^2$, le rapport entre variance inter-groupes et
variance totale) via une ANOVA à un facteur :

\begin{center}
\begin{tabular}{lcrc}
  \toprule
  \textbf{Variable} & $\boldsymbol{F}$ & $\boldsymbol{p}$ & $\boldsymbol{\eta^2}$ \\
  \midrule
  Latitude      & 29.7 & $3.1 \times 10^{-9}$ & 0.742 \\
  Longitude     & 30.0 & $2.7 \times 10^{-9}$ & 0.744 \\
  Précipitations & 10.9 & $4.6 \times 10^{-5}$ & 0.514 \\
  \bottomrule
\end{tabular}
\end{center}

L'$\eta^2$ de la latitude est de \textbf{0.742} : les 4 régions expliquent
74\% de la variance de la latitude. Pour la longitude, c'est \textbf{0.744}.
Ces deux variables combinées suffisent presque parfaitement à identifier la
région. Concrètement :
\begin{itemize}[nosep]
  \item Arctic = nord ($\text{lat} > 60°$N),
  \item Continental = centre (latitude moyenne, longitude intérieure),
  \item Atlantic = est (longitude faible, côte atlantique),
  \item Pacific = ouest (longitude élevée, côte pacifique).
\end{itemize}

\textbf{Conclusion :} les variables $Z$ sont quasi-suffisantes pour retrouver
les régions. Toute distance basée sur $D_s$ partira donc avec un avantage
structurel écrasant.

% ─────────────────────────────────────────────────────────────────────────────
\subsection{Argument 3 : rapport signal/bruit — $D_s$ sépare mieux que $D_0$}

On partage les $\binom{35}{2} = 595$ paires de stations en
\textbf{paires intra-cluster} (même région) et \textbf{paires inter-cluster}
(régions différentes). Pour une distance discriminante, on veut que les
distances inter soient nettement plus grandes que les distances intra. On
mesure le \textbf{ratio inter/intra} :

\begin{center}
\begin{tabular}{lccc}
  \toprule
  \textbf{Distance} & $\overline{d}_{\text{intra}}$ & $\overline{d}_{\text{inter}}$ & \textbf{Ratio} \\
  \midrule
  $D_0$ (fonctionnel) & 88.8  & 185.2 & 2.09 \\
  $D_s$ (vectoriel)   & 1.07  & 2.63  & 2.46 \\
  \bottomrule
\end{tabular}
\end{center}

$D_s$ a un ratio inter/intra de 2.46, contre 2.09 pour $D_0$. Cela signifie
que \textbf{$D_s$ sépare 18\% mieux les régions que $D_0$}. L'écart peut
sembler modeste, mais il est suffisant pour que la silhouette
(qui compare $a(i)$ et $b(i)$, cf.~section~\ref{sec:eval}) privilégie
systématiquement $D_s$.

\paragraph{Interprétation.} Une distance intra-cluster élevée pour $D_0$
signifie que deux stations de la même région ont des courbes
\textbf{assez différentes} (ex.~: Montréal et Halifax, toutes deux Atlantic,
ont des hivers d'intensité différente). En revanche, leur $D_s$ est faible
car elles ont des latitudes et longitudes proches. Le fonctionnel ajoute donc
de la \textbf{dispersion intra-cluster}, ce qui dégrade la silhouette.

% ─────────────────────────────────────────────────────────────────────────────
\subsection{Argument 4 : chevauchement des distributions}

On mesure le \textbf{chevauchement} entre les distributions de distances intra
et inter-cluster. On calcule le pourcentage de paires inter-cluster dont la
distance est inférieure à la médiane des distances intra-cluster (i.e.
des paires de régions différentes qui sont plus proches que la moitié des
paires de la même région) :

\begin{center}
\begin{tabular}{lc}
  \toprule
  \textbf{Distance} & \textbf{Chevauchement} \\
  \midrule
  $D_0$ (fonctionnel) & 12\% \\
  $D_s$ (vectoriel)   & 3\%  \\
  \bottomrule
\end{tabular}
\end{center}

Avec $D_0$, \textbf{12\% des paires inter-cluster sont confondues avec des
paires intra-cluster}. Avec $D_s$, ce chiffre tombe à 3\%. Cela confirme que
$D_s$ produit des clusters mieux séparés : les erreurs de frontière sont 4
fois plus rares.

% ─────────────────────────────────────────────────────────────────────────────
\subsection{Argument 5 : redondance FPCA–latitude}

Le dernier argument est peut-être le plus révélateur. On calcule les
corrélations entre les scores FPCA et les variables vectorielles :

\begin{center}
\begin{tabular}{lcc}
  \toprule
  & \textbf{Latitude} & \textbf{Longitude} \\
  \midrule
  $\xi_1$ (PC1, 88.5\%) & $-0.861$ & $-0.140$ \\
  $\xi_2$ (PC2, 8.5\%)  & $-0.184$ & $+0.100$ \\
  \bottomrule
\end{tabular}
\end{center}

La corrélation $\text{cor}(\xi_1, \text{latitude}) = -0.861$ est accablante.
Elle signifie que \textbf{le premier score FPCA est essentiellement la latitude
sous un autre nom}. La FPCA extrait le niveau moyen de température comme
premier mode de variation — et la latitude détermine ce niveau moyen.

\textbf{Conséquence pour la stratégie A.} Quand on concatène
$(\xi_1, \xi_2, \text{lat}, \text{lon}, \text{precip})$ dans $\R^5$, la
latitude apparaît \textbf{deux fois} : une fois comme $\xi_1$ et une fois
directement. Le k-means profite de ce doublement pour donner plus de poids
au niveau moyen. C'est pourquoi la stratégie A obtient le meilleur ARI
— non pas parce qu'elle combine deux sources complémentaires, mais parce
qu'elle \textbf{double implicitement} l'information géographique.

\textbf{Conséquence pour la stratégie B.} Quand on construit
$\Dw = \sqrt{\omega \cdot \Dp^2 + (1-\omega) \cdot \tilde{D}_s^2}$, ajouter
$\Dp$ (qui est essentiellement $D_0$, elle-même déterminée par $\xi_1$,
elle-même déterminée par la latitude) n'apporte \textbf{aucune information
nouvelle} par rapport à $D_s$, tout en ajoutant du bruit intra-cluster.
Le grid search conclut rationnellement : $\omega = 0$.

% ─────────────────────────────────────────────────────────────────────────────
\subsection{Synthèse du diagnostic}

\begin{center}
\fbox{\parbox{0.9\textwidth}{
\textbf{Sur Canadian Weather, les courbes de température sont une conséquence
de la géographie.} La latitude détermine le niveau moyen (= PC1 = 88.5\% de
la variance fonctionnelle), et la longitude distingue les côtes. Toute
l'information exploitable est déjà dans $Z = (\text{lat}, \text{lon},
\text{precip})$. Les courbes n'apportent pas d'information
\textbf{complémentaire} — elles répètent, en moins propre, ce que $Z$ dit
déjà.
}}
\end{center}

\subsection{Conséquence : nécessité d'un autre jeu de données}

Ce diagnostic n'invalide pas le pipeline à deux paramètres $(\alpha, \omega)$.
Il montre que \textbf{Canadian Weather n'est pas le bon terrain d'évaluation}
pour une méthode de combinaison fonctionnel–vectoriel. Pour que le grid search
trouve des valeurs $\alpha, \omega \in (0, 1)$ (pas aux coins), il faut un jeu
de données satisfaisant trois conditions :
\begin{enumerate}
  \item Les courbes portent une information \textbf{non redondante} avec $Z$ :
    la corrélation entre les scores FPCA et les variables vectorielles doit être
    faible.
  \item Les dérivées ($D_1$) ont un pouvoir discriminant propre : les
    \textbf{formes} des courbes (pas seulement les niveaux) diffèrent entre
    clusters.
  \item Les variables vectorielles \textbf{ne suffisent pas} seules : $D_s$
    seul produit un ARI modeste, de sorte que le fonctionnel peut apporter
    un gain mesurable.
\end{enumerate}

C'est ce type de jeu de données que nous cherchons pour la prochaine étape du
travail.

% ============================================================================
\section{Figures produites}

\begin{table}[H]
\centering
\begin{tabular}{lll}
  \toprule
  \textbf{Figure} & \textbf{Fichier} & \textbf{Contenu} \\
  \midrule
  Fig.~1 & \texttt{fig01\_gcv.png}             & Courbe GCV, choix de $\lambda$ \\
  Fig.~2 & \texttt{fig02\_courbes\_regions.png} & 35 courbes colorées par région \\
  Fig.~3 & \texttt{fig03\_fonctions\_propres.png} & $\varphi_1, \varphi_2$ \\
  Fig.~4 & \texttt{fig04\_scores\_fpca.png}     & Scores $\xi_1$ vs $\xi_2$ \\
  Fig.~5 & \texttt{fig05\_gridsearch\_2D.png}   & Heatmap $(\alpha, \omega)$ \\
  Fig.~5b& \texttt{fig05b\_DK\_alpha.png}       & Grid search $\alpha$ pour $\DK$ \\
  Fig.~6 & \texttt{fig06\_comparaison.png}      & Barplot comparatif \\
  Fig.~7 & \texttt{fig07\_courbes\_par\_cluster.png} & Courbes par cluster \\
  \bottomrule
\end{tabular}
\caption{Liste des figures générées.}
\end{table}

\newpage
% ============================================================================
% ============================================================================
\part{Conclusions et perspectives}
% ============================================================================
% ============================================================================

\section{Bilan de l'étude Canadian Weather}

\begin{enumerate}
  \item Le \textbf{pipeline complet} est opérationnel : du lissage au
    clustering, en passant par la FPCA et 6 distances différentes.

  \item Le système à \textbf{deux paramètres} $(\alpha, \omega)$ fonctionne
    correctement. Il permet d'explorer un espace de distances plus riche
    que le simple $\omega$ initial.

  \item Sur Canadian Weather, la meilleure stratégie est \textbf{A (FPCA + Z
    $\to$ k-means)} avec un ARI de 0.748. Les stratégies B et C ne parviennent
    pas à tirer profit de la composante fonctionnelle au-delà de ce que $Z$
    apporte seul.

  \item Ce résultat n'est \textbf{pas un échec méthodologique} : il reflète une
    propriété du jeu de données. Les variables géographiques (latitude,
    longitude) sont des prédicteurs quasi-parfaits du climat canadien. Les
    courbes de température sont une \textbf{conséquence} de la géographie,
    pas une information complémentaire.
\end{enumerate}

\section{Nécessité d'un nouveau jeu de données}
\label{sec:nouveau_dataset}

Pour valider l'intérêt du système $(\alpha, \omega)$, il faut un jeu de données
où :

\begin{itemize}
  \item La composante fonctionnelle porte une information \textbf{non
    redondante} avec les variables vectorielles.
  \item Les dérivées ($D_1$) contribuent à la discrimination (les formes
    importent, pas seulement les niveaux).
  \item Les variables vectorielles \textbf{ne suffisent pas} à elles seules
    pour un bon clustering.
\end{itemize}

Un tel jeu de données permettrait de vérifier que le grid search 2D trouve des
valeurs optimales $\alpha, \omega \in (0, 1)$ (pas aux coins), prouvant ainsi
l'utilité de la combinaison.

\section{Reproductibilité}

\begin{table}[H]
\centering
\begin{tabular}{ll}
  \toprule
  \textbf{Paramètre} & \textbf{Valeur} \\
  \midrule
  Seed       & 42 \\
  R version  & $\geq 4.0$ \\
  Packages   & \texttt{fda}, \texttt{cluster}, \texttt{mclust} \\
  nbasis     & 65 \\
  norder     & 4 (B-splines cubiques) \\
  $\lambda$  & 8.15 (GCV) \\
  $K$        & 2 (seuil 95\%) \\
  $k$        & 4 clusters \\
  Grille $\alpha$ & $\{0, 0.1, \ldots, 1\}$ \\
  Grille $\omega$ & $\{0, 0.1, \ldots, 1\}$ \\
  \bottomrule
\end{tabular}
\caption{Tous les paramètres du pipeline.}
\end{table}

\bigskip

L'intégralité du code source est dans \texttt{src/}. Pour reproduire les résultats :
\begin{lstlisting}
Rscript src/main.R
\end{lstlisting}

% ============================================================================
\begin{thebibliography}{9}

\bibitem{ramsay2005}
  J.O.~Ramsay et B.W.~Silverman,
  \textit{Functional Data Analysis}, 2\textsuperscript{e} édition,
  Springer, 2005.

\bibitem{jaakkola1999}
  T.~Jaakkola et D.~Haussler,
  \textit{Exploiting Generative Models in Discriminative Classifiers},
  NIPS, 1999.

\bibitem{ferreira2014}
  L.~Ferreira et B.~de Carvalho,
  \textit{Kernel-based Clustering of Functional Data},
  2014.

\end{thebibliography}

\end{document}
